% Source: https://artofproblemsolving.com/community/c3h1225611_hard_question
% Stable identity: aops-topic-1225611-q01
Let the squares have side $1$ and put $E=(0,1)$, $F=(1,1)$, $B=(1,0)$, $C=(2,0)$, $D=(3,0)$. Measured from the ray $EF$, the segments $EB$, $EC$, $ED$ go $1$ down for $1$, $2$, $3$ across, so $\tan\angle BEF=1$, $\tan\angle CEF=\dfrac{1}{2}$, $\tan\angle DEF=\dfrac{1}{3}$. Thus $\angle BEF=45^\circ$.

\begin{asy}
usepackage("amsmath");
size(6cm);
defaultpen(fontsize(10pt));
draw((0,0)--(3,0)--(3,1)--(0,1)--cycle,linewidth(0.9)); draw((1,0)--(1,1)); draw((2,0)--(2,1));
draw((0,1)--(1,0),blue); draw((0,1)--(2,0),red); draw((0,1)--(3,0),heavygreen); label("$A$",(0,0),SW); label("$B$",(1,0),S); label("$C$",(2,0),S); label("$D$",(3,0),SE); label("$E$",(0,1),NW); label("$F$",(1,1),N); label("$G$",(2,1),N); label("$H$",(3,1),NE);
label("\mbox{}",(min(currentpicture,user=true)+max(currentpicture,user=true))/2,black); // keeps the text after the figure black when LaTeX typesets the labels inline
\end{asy}

For the other two, use the tangent addition formula:
\[
\tan(\angle CEF+\angle DEF)=\dfrac{\dfrac{1}{2}+\dfrac{1}{3}}{1-\dfrac{1}{2}\cdot\dfrac{1}{3}}=\dfrac{\dfrac{5}{6}}{\dfrac{5}{6}}=1,
\]
and both angles are acute with sum less than $90^\circ$, so $\angle CEF+\angle DEF=45^\circ$. (Equivalently, with complex numbers: the directions of $EC$ and $ED$ are $2-i$ and $3-i$, and $(2-i)(3-i)=5-5i$ has argument $-45^\circ$, so the two angles add to $45^\circ$.) Hence the three angles add to $45^\circ+45^\circ=90^\circ$. $\blacksquare$
